现考虑 $\tau^{-1}$ 线性算子

\begin{equation}\tag{4.5.13}
\tau^{-1}\circ\psi_p\circ H(X):L(b)\longrightarrow L(b).
\end{equation}

其迭代为

\begin{equation}\tag{4.5.14}
(\tau^{-1}\circ\psi_p\circ H(X))^a
=\tau^{-a}\circ\psi_{p^a}\circ
H(X)H^\tau(X^p)\cdots H^{\tau^{a-1}}(X^{p^{a-1}}),
\end{equation}

且若 $\operatorname{ord}(q)$ 整除 $a$, 则 $\tau^{-a}=\operatorname{id}$.
定义 $\tau^{-1}$ 线性运算

\begin{equation}\tag{4.5.15}
\left\{
\begin{aligned}
\alpha:L^0\!\left(\frac1{p-1}\right)
&\longrightarrow L^0\!\left(\frac1{p-1}\right),\\
\alpha&=\frac1p\,\tau^{-1}\circ\psi_p\circ H(X).
\end{aligned}
\right.
\end{equation}

则 (4.5.12) 具有形式

\begin{equation}\tag{4.5.16}
\operatorname{card}V^*(\mathbb F_{q^\nu})
=(q^\nu-1)^{n+1}
+(q^\nu-1)^{n+2}
\operatorname{tr}\!\left(
\left.\alpha^{\operatorname{ord}(q^\nu)}
\right|L^0\!\left(\frac1{p-1}\right)\right).
\end{equation}

置

\begin{equation}\tag{4.5.17}
S=\{1,\ldots,n+2\}.
\end{equation}

对每个非空子集 $A\subset S$, 定义子空间

\begin{equation}\tag{4.5.18}
\begin{aligned}
L^A(b)=\{\,&\text{级数 $\sum_U B_UX^U\in L(b)$, 满足对每个 $B_U\ne0$ 的单项式 $X^U$,}\\
&\text{使 $U_i>0$ 的指标 $i$ 恰为 $0$ 以及 $j\in A$}\,\}.
\end{aligned}
\end{equation}

则

\begin{equation}\tag{4.5.19}
L^0(b)=\bigoplus_{\substack{A\subset S\\A\ne\varnothing}}L^A(b).
\end{equation}

\medskip
\noindent\emph{底本注.}
印本在 (4.5.10) 的因子 $q^\nu-1$ 以及 (4.5.12) 的分母 $q^\nu$ 中漏去 $\nu$,
尽管这两个公式均直接改写自 (4.4.6). 它还将公式编号 (4.5.14) 颠倒为 (5.4.14).
照录法文层保留这三处印本形式; 校订法文层和英文层采用上列由代数关系强制确定的读法.
